% Plain-words narration of the paper -- a standalone companion document.
% Build on its own:  pdflatex narration
% This is NOT part of the submission (main.tex); it explains the paper section by
% section in simple words, showing how each part leads into the next.
\documentclass[11pt]{article}
\usepackage[margin=1in]{geometry}
\usepackage{parskip}
\usepackage{enumitem}
\usepackage{titlesec}
\titleformat{\section}{\large\bfseries}{\thesection.}{0.5em}{}
\titleformat{\subsection}{\normalsize\bfseries}{}{0em}{}
\setlist{nosep,leftmargin=1.4em}
\newcommand{\handoff}[1]{\textit{Hand-off.}\ #1}
\newcommand{\says}[1]{\textbf{What it says.}\ #1}
\newcommand{\why}[1]{\textbf{Why it's here.}\ #1}

\title{\vspace{-2em}Reading the paper as one story\\[-0.3em]
  {\normalsize\normalfont (plain words)}}
\author{}
\date{}

\begin{document}
\maketitle
\vspace{-3em}

This walks the paper from top to bottom. For each part it says, in simple words, what the
part says, why it is there, and how it hands off to the next part---so the whole thing
reads as one connected story, not a pile of sections. No jargon.

\section*{The whole story in five sentences}
\begin{enumerate}
\item AI agents now write parallel code and check their own work with a test.
\item That same test is what new models are \emph{trained} on, so the test decides what
  ``good code'' means.
\item For parallel code the usual test---``is the answer correct?''---is the wrong test,
  because plain slow code passes it perfectly.
\item So a model trained to pass it can quietly learn to write correct code that is not
  actually fast.
\item The fix is to test two things, correct \emph{and} fast, and we show what that
  changes.
\end{enumerate}
Everything below is just this story, told carefully.

\section{Abstract --- the promise}
\says{In one paragraph: agents check their own parallel code; that check is also the
training reward; the correctness-only check is both \emph{already solved} (so a fancier
one adds nothing) and \emph{cheatable} (plain serial code passes it); therefore the check
must also measure speed; and here is the evidence.}

\why{It is the trailer. A reviewer decides in 30 seconds whether to care.}

\handoff{It names two surprises---``correctness is solved'' and ``correct code is not
always fast.'' The introduction now has to make you \emph{believe} those two things.}

\section{Introduction --- the problem, with a picture}
\says{Agents write parallel code. Parallel code is special: whether the answer is right
can depend on \emph{how many threads run it}. Then the picture that carries the whole
paper---the \textbf{histogram}: an agent counts items into bins in parallel, two threads
bump the same counter at once, one count is lost. On \emph{one} thread there is no
collision, so it looks perfect. The agent ran it once, saw a good-looking answer, and
shipped a bug. The check the agent did was blind to the bug.}

\why{It plants one idea you never forget: the cheap check misses exactly the bugs parallel
code has. Then it names the key move---the check an agent uses is the same signal a model
is trained on, so \textbf{the check is the definition of good code}.}

\handoff{If the check is that important, we had better define it exactly. That is the next
part.}

\section{The metric --- defining the check precisely}
\says{It gives the check a name: \textbf{configuration-robust correctness}. Plain version:
do not just run the code once---run it at 1, 2, 4, and 8 threads, several times each, and
compare every run to a trusted simple (serial) answer. Only call it correct if
\emph{every} run matches. It also says the quiet thing out loud: this check looks only at
the \emph{answer}, so a program with all the parallel parts deleted---plain slow
code---passes it every time. And it introduces the second measuring stick,
\textbf{speed} (how much faster 8 threads are than 1).}

\why{It turns the fuzzy word ``correct'' into something you can measure and argue about,
and it plants the seed of the twist: a correct-only check is cheatable.}

\handoff{Now we have a precise check. Time to actually use it on real models. That is the
setup.}

\section{Setup --- what we actually tested}
\says{The benchmark is \textbf{ParEval} (standard tasks for parallel code, each with a
trusted simple answer we compare against). We test a handful of strong models. We build
the check as a real tool the agent can call, at three strengths: no tool (just generate),
a \emph{weak} tool (test at 1 thread only), and the \emph{full} tool (test across all
thread counts). Same agent, same prompts---only the tool changes. That controlled swap is
how we isolate what the strong check is worth.}

\why{It makes the results fair and repeatable: any difference later must come from the
tool, not from luck.}

\handoff{Everything is set. What happened? The results.}

\section{Results --- two surprises}
\says{Two findings, in order.}
\begin{itemize}
\item \textbf{Surprise one: the strong check changes nothing.} For these good models,
  testing at all thread counts gives the exact same result as testing at one thread. Why?
  They simply do not write the sneaky race the strong check is built to catch---and an
  independent race detector agrees the accepted code is clean. So correctness is
  \emph{already solved} here, and a finer correctness check has nothing left to catch.
  (This is the honest ``negative result.'')
\item \textbf{Surprise two: correct does not mean fast.} We take every program the check
  accepted and actually time it. They range from \textbf{7.9$\times$ faster on 8 threads
  down to 0.33$\times$}---one accepted program runs \emph{slower} on 8 threads than on 1.
  The check rated them all equally correct. So the check cannot tell a great parallel
  program from a fake one.
\end{itemize}

\why{Surprise one clears correctness off the table. Surprise two shows the real problem is
now speed. Together they force a question.}

\handoff{If a correct-only check is both pointless-at-the-top and cheatable, then
\textbf{what should the check---and the training reward---measure instead?} That question
is the next section, and it is the heart of the paper.}

\section{The fix --- a two-part reward, and four experiments}
\says{Make the reward measure \textbf{two things}: first it must be correct, and only then
does it score how fast it is. Plain slow code clears the ``correct'' bar but scores near
the bottom on speed; correct-and-fast code scores near the top. Then we test this idea
four ways, each one a step more committed:}
\begin{enumerate}
\item \textbf{What the old reward throws away (done, real numbers).} Among tasks where
  several different correct programs exist, the correct-only reward treats them all the
  same, so picking one is a coin flip---you get the \emph{average} speed. The two-part
  reward picks the \emph{fastest}. The gap is real: average $\sim$3.5$\times$ vs
  $\sim$4.2$\times$, and on the worst task it is the difference between a 0.33$\times$ and
  a 6.8$\times$ program the old reward calls equal.
\item \textbf{Best-of-N with live models.} Generate several answers, score them, keep the
  best. Under the old reward you keep a random correct one; under the two-part reward you
  keep the fast one. This is a cheap stand-in for what RL training would learn.
\item \textbf{The two-gate agent.} A real agent pipeline: first a ``coder'' gets it
  correct (it calls the checker itself), then an ``optimizer'' makes it fast (it calls the
  timer itself)---and every speed change is re-checked for correctness, so speed can never
  sneak in a wrong answer.
\item \textbf{Actual training (GRPO).} Train a model with each reward and watch. Predict:
  old reward $\rightarrow$ the code slowly gets less parallel; two-part reward
  $\rightarrow$ it stays fast.
\end{enumerate}

\why{It turns the complaint (``the reward is wrong'') into a fix (``here is the better
reward'') and backs it with evidence at four levels, from a paper-only analysis up to real
training.}

\handoff{Now step back: what does all this mean, and where could we be wrong?}

\section{Discussion --- meaning and honesty}
\says{Meaning: the hard part of parallel-code writing moved. It used to be ``do not write
races''; strong models mostly cleared that, so now it is ``actually be fast,'' and a
correct-only reward is blind to that. Honesty: we tested shared-memory OpenMP up to 8
threads, a handful of models, a limited task set---so the exact numbers may not carry to
other settings; but the \emph{speed} finding is directly measured, not inferred, so it is
solid.}

\why{It tells the reader how far to trust the result, which is what good reviewers look
for.}

\handoff{Before claiming it is new, place it next to what others did.}

\section{Related work --- who did what, and why we are different}
\says{Rewarding code models for \emph{speed} is \textbf{not new}---several groups do it,
including the group that made ParEval. So we do \textbf{not} claim ``add speed to the
reward.'' Our different, honest claim is the \textbf{diagnosis}: a correct-only reward is
cheatable and saturated, we \emph{measure} how much speed it throws away, and we add a
speed gate that is guarded by correctness. We say this plainly so a reviewer cannot say
``this has been done.''}

\why{It protects the paper. Naming the neighbors and drawing a clean line is what keeps it
from being desk-rejected as ``already known.''}

\handoff{One last tightening of the message.}

\section{Conclusion --- the one thing to remember}
\says{Correctness stopped being the hard part for strong models; ``correct but not fast''
is where their parallel code fails now, and any check---or training reward---that looks
only at the answer cannot see it. Measure both.}

\why{It leaves the reader with the single sentence you want repeated back.}

\section*{The thread, in one line each}
\begin{itemize}
\item \textbf{Abstract:} the check is the reward, and the correct-only check is the wrong
  one.
\item \textbf{Intro:} here is why parallel code fools a one-run check (the histogram).
\item \textbf{Metric:} here is the check, stated exactly---and yes, plain slow code passes
  it.
\item \textbf{Setup:} here is how we test it fairly (swap only the tool).
\item \textbf{Results:} the strong check adds nothing, and correct code ranges from fast
  to slower-than-serial.
\item \textbf{Reward:} so measure both---and here are four experiments showing it matters.
\item \textbf{Discussion:} the hard part moved from correctness to speed; here is what we
  can and cannot claim.
\item \textbf{Related work:} speed-rewards exist; our new part is the diagnosis and the
  guarded speed gate.
\item \textbf{Conclusion:} measure both axes, or you are grading the wrong thing.
\end{itemize}

Each arrow is a \emph{because}: because the check is the reward $\rightarrow$ we define it;
because it is cheatable $\rightarrow$ we test what that costs; because that costs real
speed $\rightarrow$ we add a speed gate; because speed-rewards exist $\rightarrow$ we are
careful to claim only the diagnosis. That is the story.

\end{document}
